\begingroup
\makeatletter\tagsleft@true\let\veqno\@@leqno\makeatother
\clearpage\SourceLeaf{323b}\ArithmeticHead{316}{L}{ANALYSIS.\quad NACHLASS.}
\NumberedItem{6}
\par\indent Ist [log] $t$ die Differenz, wenn man von dem vollständigen Potential (d. i. imaginären Theil mit gerechnet) der im Centrum concentrirten Masse das vollständige Potential der auf die Ellipse vertheilten Masse subtrahirt, beide in dem Punkte geltend, dessen complexe Zahl $u+\frac\lambda u$ ist, so wird
\clearpage\SourceLeaf{324}\ArithmeticHead{317}{R}{{\scriptsize ABBILDUNG DER ELLIPSE AUF DEN KREIS.}}
\begin{equation}
\begin{aligned}
\log t={}&\log u-\frac\lambda{1+\lambda\lambda}(uu-u^{-2})+\frac12\frac{\lambda\lambda}{1+\lambda^4}(u^4-u^{-4})\\
&-\frac13\frac{\lambda^3}{1+\lambda^6}(u^6-u^{-6})\quad\text{u.s.w.,}
\end{aligned}\tag*{[14]}
\end{equation}
\noindent welches die obige Reihe selbst ist.\quad[Siehe die Gl. [7], oben S. 314].
\par\indent Hieraus ferner
\begin{equation}
t=u.\frac{1+\lambda u^{-2}}{1+\lambda uu}.\frac{1+\lambda^3uu}{1+\lambda^3u^{-2}}.\frac{1+\lambda^5u^{-2}}{1+\lambda^5uu}.\text{etc.}\tag*{[15]}
\end{equation}
\par\indent Man kann $t$ auch folgende Form geben, wobei
\[\mu\mu=\lambda\]
\noindent gesetzt ist,
{\small\begin{equation}
\begin{aligned}
t&=\frac{1+\mu^4\left(\dfrac1{\mu\mu uu}+\mu\mu uu\right)+\mu^{16}\left(\dfrac1{\mu^4u^4}+\mu^4u^4\right)+{.}{.}{.}}{\mu\left(\dfrac1{\mu u}+\mu u\right)+\mu^9\left(\dfrac1{\mu^3u^3}+\mu^3u^3\right)+\mu^{25}\left(\dfrac1{\mu^5u^5}+\mu^5u^5\right)+{.}{.}{.}}\\[7pt]
&=u\frac{1+\mu\mu u^{-2}+\mu^6u^2+\mu^{12}u^{-4}+\mu^{20}u^4+{.}{.}{.}}{1+\mu\mu u^2+\mu^6u^{-2}+\mu^{12}u^4+\mu^{20}u^{-4}+{.}{.}{.}}.
\end{aligned}\tag*{[16]}
\end{equation}}
\noindent {\small[S. 224]}
\par\indent Noch zierlicher setze man
\[u=\lambda^\alpha.\]
\noindent Es ist dann
\begin{equation}
t=\frac{\sum\lambda^{\frac18(4n+1+2\alpha)^2}}{\sum\lambda^{\frac18(4n-1+2\alpha)^2}},\tag*{[17]}
\end{equation}
\noindent wo für $n$ alle ganzen Zahlen positiv, negativ und $0$ eingeschlossen zu setzen sind.
\par\indent Daraus folgt,
\[t=f(\lambda,u)\]
\noindent gesetzt,
\[f(\lambda,\lambda u)f(\lambda,u)=1.\]
\noindent Die allgemeine Gleichung ist
\begin{equation}
\frac{dt}{\surd(aa+bbtt\,\surd(bb-aatt)}=\frac{i\,du}{u}=\frac{dT}{\surd(4\lambda-TT)},\tag*{[18]}
\end{equation}
\noindent wo $h$ mit $\lambda$ oder $\mu$ durch die Gleichung zusammenhängt
\clearpage\SourceLeaf{325}\ArithmeticHead{318}{L}{ANALYSIS.\quad NACHLASS.}
{\small\begin{equation}
h=\left(\frac{2\mu+2\mu^9+2\mu^{25}+{.}{.}{.}}{1+2\mu^4+2\mu^{16}+{.}{.}{.}}\right)^2=\frac{bb}{aa},\quad\text{\textit{bedarf Berichtigung}[*)]}\tag*{[19]}
\end{equation}}
\noindent [und] wo
\begin{equation}
ab=2(\mu^{\frac12}+\mu^{\frac92}+\mu^{\frac{25}2}+{.}{.}{.})^2.\tag*{[20]}
\end{equation}
\par\bigskip\noindent\rule{.25\linewidth}{.4pt}\par\smallskip
{\small\par\indent [*) Siehe art. [8.], unten S. 319.]\par}
\clearpage
\endgroup
